% Solutions to Chapter 2, from lectures/L02/appendix/solutions: the README's list of programs,
% and every source file typeset straight from the repository.

\section{Chapter 2: Classes}
\label{sol:sec:c2}
Reference implementations for the exercises in \secref{c2:sec:exercises}, one program per exercise
set, in \file{lectures/L02/appendix/solutions}. Each directory has a Makefile; \code{make} in the
directory builds the program and runs it.

% ------------------------------------------------------------------------------------------
\subsection{Exercise Set 1: First Class}
\label{sol:set:c2.1}
Exercises~\ref{c2:ex:1.1} and~\ref{c2:ex:1.2}, simple GPIO classes, in
\file{lectures/L02/appendix/solutions/exercise1}. Both classes are implemented directly in their
headers, and \code{main()} exercises each in turn.

\cppfile{lectures/L02/appendix/solutions/exercise1/include/driver/gpio/led.hpp}
\cppfile{lectures/L02/appendix/solutions/exercise1/include/driver/gpio/button.hpp}
\cppfile{lectures/L02/appendix/solutions/exercise1/main.cpp}

% ------------------------------------------------------------------------------------------
\subsection{Exercise Set 2: Interaction Between Classes}
\label{sol:set:c2.2}
Exercises~\ref{c2:ex:2.1} and~\ref{c2:ex:2.2}, interaction between GPIO classes, in
\file{lectures/L02/appendix/solutions/exercise2}. The two headers, \file{led.hpp} and
\file{button.hpp}, are identical to those of Exercise Set~1 above; only the program differs.

\cppfile{lectures/L02/appendix/solutions/exercise2/main.cpp}

% ------------------------------------------------------------------------------------------
\subsection{Exercise Set 3: Class Split Across Multiple Files}
\label{sol:set:c2.3}
Exercise~\ref{c2:ex:3.1}, the buzzer class, in \file{lectures/L02/appendix/solutions/exercise3}. The
declarations are in the header and the definitions in the source file, and the Makefile is the one
from \secref{c1:sec:tips}: the \file{include} directory on the compiler's search path, and both
source files listed.

\makefile{lectures/L02/appendix/solutions/exercise3/Makefile}
\cppfile{lectures/L02/appendix/solutions/exercise3/include/driver/buzzer.hpp}
\cppfile{lectures/L02/appendix/solutions/exercise3/source/driver/buzzer.cpp}
\cppfile{lectures/L02/appendix/solutions/exercise3/source/main.cpp}

% ------------------------------------------------------------------------------------------
\subsection{Exercise Set 4: Timer Class}
\label{sol:set:c2.4}
Exercise~\ref{c2:ex:4.1}, the timer class, in \file{lectures/L02/appendix/solutions/exercise4}. The
loop runs for three timeouts' worth of ticks, so with a 1000\,ms timeout it runs the 3000 iterations
the exercise asks for.

\cppfile{lectures/L02/appendix/solutions/exercise4/include/driver/timer.hpp}
\cppfile{lectures/L02/appendix/solutions/exercise4/source/driver/timer.cpp}
\cppfile{lectures/L02/appendix/solutions/exercise4/source/main.cpp}
